\section{Passive multi-port electromagnetic macromodel}
\label{sec:port-macro}

\subsection{Narrowband port operator}

For a fixed carrier frequency, the electromagnetic subsystem is summarized by
the reciprocal impedance matrix
\[
  v=Z(\mathcal S,\omega,T)i,\qquad Z^\T=Z .
\]
The real Hermitian part is decomposed into physically named channels:
\[
  \Herm Z
  =D_{\mathrm{cond}}
   +D_{\mathrm{diel}}
   +D_{\mathrm{rad}}
   +D_{\mathrm{other}},
\]
with each $D_\alpha\PSD$ for a passive scene. This decomposition is preferable
to obtaining a small outward or dielectric term as the subtraction of two large
independently approximated matrices.

\subsection{Broadband positive-real realization}

For broadband drive, modulation, or electrical transients, the port response is
represented by a rational passive state-space model
\[
 \dot x_e=A_e x_e+B_e i,\qquad
 v=C_e x_e+D_e i.
\]
A general linear port-Hamiltonian descriptor realization may be written
\cite{cherifi2024phpassive}
\[
 E_e\dot x_e
 =
 (J_e-R_e)Q_e x_e+(G_e-P_e)i,
\]
\[
 v
 =
 (G_e+P_e)\adj Q_e x_e+(S_e+N_e)i,
\]
with
\[
 Q_e\adj E_e=E_e\adj Q_e\PSD,
\]
a skew-adjoint interconnection block and a positive-semidefinite dissipation
block. In particular,
\[
 \begin{bmatrix}
 J_e&G_e\\
 -G_e\adj&N_e
 \end{bmatrix}
 =
 -
 \begin{bmatrix}
 J_e&G_e\\
 -G_e\adj&N_e
 \end{bmatrix}\adj,
\]
and
\[
 \begin{bmatrix}
 Q_e\adj R_eQ_e&Q_e\adj P_e\\
 P_e\adj Q_e&S_e
 \end{bmatrix}
 \PSD.
\]
This structure implies passivity. Positive-real transfer behavior and
port-Hamiltonian realizability are closely related but are not treated as
unconditionally identical for arbitrary descriptor systems
\cite{cherifi2024phpassive}.

For the first broadband implementation, vNext deliberately uses the simpler
collocated standard-state subclass
\[
 E_e=I,\qquad P_e=0,
\]
\[
 \dot x_e=(J_e-R_e)Q_e x_e+B_e i,
\]
\[
 v=B_e\adj Q_e x_e+(S_e+N_e)i,
\]
with
\[
 J_e=-J_e\adj,\quad N_e=-N_e\adj,\quad
 Q_e\SPD,\quad R_e\PSD,\quad S_e\PSD.
\]
The corresponding transfer matrix is
\[
 Z(s)
 =
 B_e\adj Q_e
 \big(sI-(J_e-R_e)Q_e\big)^{-1}B_e
 +S_e+N_e.
\]
Reciprocity imposes additional realization constraints so that
$Z(s)^\T=Z(s)$; the exact reciprocal subclass is fixed together with the
broadband backend rather than assumed from passivity alone.

\subsection{Structure-preserving frequency fitting}

REFERENCE samples over frequency are compressed into a rational model using a
stable fitting method such as vector fitting followed by passivity-preserving
realization or direct structure-preserving interpolation
\cite{gustavsen1999vector,benner2015mor}. Neural networks do not predict
unconstrained pole/residue tables. They predict Cholesky/factor coordinates of
the structured realization or a correction to an analytic baseline.

\subsection{Circuit coupling}

External circuits are first-class components. Let the network equations be
\[
  F_{\mathrm{circ}}(v,i,x_c,u,t)=0.
\]
The electromagnetic port model is coupled to these equations without
reconstructing spatial electromagnetic unknowns. At narrowband, one solves
\[
  v=Z(\mathcal S,\omega,T)i
\]
with the circuit algebraic constraints. In broadband mode, the electrical
macromodel state $x_e$ is integrated jointly with circuit states $x_c$ and
thermal states.

\subsection{Temperature dependence}

The model may use a separated representation
\[
 Z(s;\mathcal S,T)
 =
 Z_{\mathrm{base}}(s;\mathcal S)
 +\sum_{r=1}^{R_T}\gamma_r(T)\Delta Z_r(s;\mathcal S),
\]
or may parameterize the passive realization matrices as smooth functions of a
small thermal state. Any interpolation in temperature must preserve the
positive-real constraints.
